% Source-aware Rethlas proof context for thm:ok-mvob-line-bundles.
% Packet metadata only; component bodies are exact live-source extracts.
% Ch55 live SHA-256: b2134d6b59f7687f8f5f296023a7d80703122210909032856b965f12f61619ab
% Ch51 live SHA-256: 60e1165fad0cf2b05e5bb03602c485a18d49160b40d1fc4c2569e4ca383371dc

% BEGIN EXACT COMPONENT CH55_ALGEBRAIC_VOLUME_ORDER
% Source: manuscript/parts/part5/ch55_okounkov_dh.tex lines 97--155
% SHA-256: 55ed9b814701efe92bf5745865860555a43155d6476e660e96f06935155fe668
\begin{theorem}\label{thm:ok-algebraic-volume-order}
Let \(\varphi,\psi\in\PSH(X,\theta)\) have positive multiplier-ideal
volume.
Then
\[
 \lambda_n\!\left(\Delta_\nu(\theta,\varphi)\right)
     =\frac1{n!}\vol(\theta,\varphi).
\tag{55.4}\label{eq:ok-algebraic-volume}
\]
If
\(\varphi\preceq_{\cI}\psi\), meaning
\(\cI(k\varphi)\subseteq\cI(k\psi)\) for every \(k>0\), then
\[
 \Delta_\nu(\theta,\varphi)
       \subseteq\Delta_\nu(\theta,\psi).
\tag{55.5}\label{eq:ok-i-order}
\]
In particular the body depends only on the \(\cI\)-singularity type.
For a potential with minimal singularities it is the usual Okounkov body
\(\Delta_\nu(L)\).
\end{theorem}

\begin{proof}[Almost-semigroup source boundary]
The ideal inclusion gives
\(W_k(\theta,\varphi)\subseteq W_k(\theta,\psi)\) for every \(k\).
Monotonicity of the almost-semigroup body gives
\eqref{eq:ok-i-order}.

The remaining assertion is the central almost-semigroup theorem of
\cite{Xia21}, reproduced in the supplied Okounkov chapter of
\cite{XiaBook}.  We record its inputs so that the boundary is visible.
First, the multiplier-ideal family
\(\Gamma_\nu(\theta,\varphi)\) is proved to be an almost semigroup, and its
asymptotic counting volume is
\[
 \lim_{k\to\infty}k^{-n}
 h^0\!\left(X,L^k\otimes\mathcal I(k\varphi)\right)
 =\frac1{n!}\vol(\theta,\varphi).
\]
This is the multiplier-ideal asymptotic Riemann--Roch theorem.  Second, the
almost-semigroup body construction is continuous for its pseudometric and
has Euclidean volume equal to the counting volume.

For orientation, the source proves these inputs first for analytic
singularities defining a K\"ahler current.  On a log resolution
\[
                   \pi^*\theta_\varphi=[D]+R,
\tag{55.6}\label{eq:ok-log-resolution}
\]
divisorial valuations translate the ordinary Okounkov body by
\(\nu(D)\), and asymptotic Riemann--Roch computes the residual volume.  It
then passes to K\"ahler currents by quasi-equisingular approximation and
to arbitrary positive-volume potentials by
\(\varphi_\varepsilon=(1-\varepsilon)\varphi+\varepsilon\chi\).
Continuity is required both for the almost semigroups and for
\(\mathcal I\)-model volume.  Applying the two imported inputs to
\eqref{eq:ok-partial-body} gives
\eqref{eq:ok-algebraic-volume}.
\end{proof}
% END EXACT COMPONENT CH55_ALGEBRAIC_VOLUME_ORDER

% BEGIN EXACT COMPONENT CH55_I_MODEL_MASS_VOLUME
% Source: manuscript/parts/part5/ch55_okounkov_dh.tex lines 370--390
% SHA-256: 17d9f654213fafafc73940a6a0cd5fe1841c8dae81c0d357ebf68852b63c1c64
\begin{lemma}\label{lem:ok-i-model-mass-volume}
If \(u\in\PSH(X,\theta)\) is \(\mathcal I\)-model, then
\[
                         \vol(\theta,u)=M_\theta(u).
\]
In particular, every interior level of a positive-mass
\(\mathcal I\)-model test curve lies in the domain of
\Cref{def:ok-algebraic-partial-body}.
\end{lemma}

\begin{proof}
The defining equality \(P_\theta[u]_{\mathcal I}=u\), inserted into
\eqref{eq:complex-current-volume}, gives
\[
 \vol(\theta,u)
 =\int_X\bigl(\theta+\ddc u\bigr)^n
 =M_\theta(u).
\]
The last assertion follows from the positive-mass condition on the interior
levels in \Cref{def:rays-test-curve}.
\end{proof}
% END EXACT COMPONENT CH55_I_MODEL_MASS_VOLUME

% BEGIN EXACT COMPONENT CH51_I_ENVELOPE_PROJECTION
% Source: manuscript/parts/part5/ch51_complex_prerequisites.tex lines 3240--3288
% SHA-256: 54e4d6b864e81429549753d781a3e1c859cb5235a0d138b556196214611dbaa0
\begin{definition}[\(\mathcal I\)-envelope]
\label{def:complex-i-envelope}
For \(\varphi\in\PSH(X,\theta)\), set
\[
 P_\theta[\varphi]_{\mathcal I}
 :=\uscsup\{\psi\in\PSH(X,\theta)\mid
          \psi\leq0,\ \psi\sim_{\mathcal I}\varphi\}.
\tag{51.27}\label{eq:complex-i-envelope}
\]
It is \(\mathcal I\)-model if it is fixed by this operator.
\end{definition}

\begin{theorem}\label{thm:complex-i-projection}
For every \(\varphi\in\PSH(X,\theta)\),
\[
 \varphi\sim_{\mathcal I}P_\theta[\varphi]_{\mathcal I},
\qquad
 P_\theta[P_\theta[\varphi]_{\mathcal I}]_{\mathcal I}
      =P_\theta[\varphi]_{\mathcal I}.
\tag{51.28}\label{eq:complex-i-projection}
\]
Furthermore,
\[
 P_\theta[\varphi]\leq P_\theta[\varphi]_{\mathcal I},
\qquad
 \varphi\sim_{\mathcal I}P_\theta[\varphi].
\tag{51.29}\label{eq:complex-p-i-comparison}
\]
If \(\varphi\) has analytic singularities, then
\[
 \varphi\sim P_\theta[\varphi]
        \sim P_\theta[\varphi]_{\mathcal I}.
\tag{51.30}\label{eq:complex-analytic-envelope-agreement}
\]
\end{theorem}

\begin{proof}
Choquet reduction replaces the family in
\eqref{eq:complex-i-envelope} by an increasing sequence.  Finite maxima
preserve divisorial Lelong numbers, and strong openness shows that its
usc supremum has the same multiplier ideals as \(\varphi\).  This proves
\eqref{eq:complex-i-projection}.  Formula
\eqref{eq:complex-p-envelope-rooftop}, again with strong openness, gives
\(\varphi\sim_{\mathcal I}P_\theta[\varphi]\), hence
\eqref{eq:complex-p-i-comparison}.  For analytic singularities, pass to
a log resolution.  Divisorial orders then control the whole singularity
up to \(O(1)\), so the two envelopes cannot be strictly less singular
than \(\varphi\).
\end{proof}
% END EXACT COMPONENT CH51_I_ENVELOPE_PROJECTION

% BEGIN EXACT COMPONENT CH51_CURRENT_VOLUME
% Source: manuscript/parts/part5/ch51_complex_prerequisites.tex lines 3488--3532
% SHA-256: 26c495f9d8ec7d043106ceb6b558205407202047e902c28a113b9c8269db917b
\subsection{The volume detected by ideals}
\label{sec:ch51-current-volume}

\begin{definition}[Volume detected by ideals]
\label{def:complex-current-volume}
For \(\varphi\in\PSH(X,\theta)\), define
\[
 \vol(\theta,\varphi)
 :=\int_X\bigl(\theta+\ddc
             P_\theta[\varphi]_{\mathcal I}\bigr)^n.
\tag{51.32}\label{eq:complex-current-volume}
\]
It depends only on the current \(\theta_\varphi\).  Moreover,
\[
 M_\theta(\varphi)\leq\vol(\theta,\varphi),
\tag{51.33}\label{eq:complex-mass-volume}
\]
with equality for analytic singularities.
\end{definition}

Three comments.  The nonpluripolar product on the right of
\eqref{eq:complex-current-volume} is, as always, the one of
\eqref{eq:complex-global-npp}; the notation omits the angle brackets only
because the argument is now a specific model potential.  The inequality
\eqref{eq:complex-mass-volume} is
\Cref{thm:complex-npp-ledger} together with
\(P_\theta[\varphi]\leq P_\theta[\varphi]_{\mathcal I}\) from
\eqref{eq:complex-p-i-comparison} and the mass identity
\eqref{eq:complex-p-projection}; explicitly,
\[
 M_\theta(\varphi)
 =M_\theta\bigl(P_\theta[\varphi]\bigr)
 \leq M_\theta\bigl(P_\theta[\varphi]_{\mathcal I}\bigr)
 =\vol(\theta,\varphi),
\]
the middle step because \(P_\theta[\varphi]\preceq
P_\theta[\varphi]_{\mathcal I}\).  Equality for analytic singularities is
\eqref{eq:complex-analytic-envelope-agreement}.  Finally, the dependence
only on the current is the statement that
\(\vol(\theta,\varphi)\) is unchanged when \(\varphi\) is replaced by
\(\varphi+c\), which is
\Cref{prop:ch51-mass-inequalities}, item
\cref{itm:ch51-mi:translation}, together with the fact that
\eqref{eq:complex-i-envelope} depends only on the
\(\mathcal I\)-class, which is invariant under adding constants.
% END EXACT COMPONENT CH51_CURRENT_VOLUME

% BEGIN EXACT COMPONENT CH55_METRIC_ONLY_CHAIN
% Source: manuscript/parts/part5/ch55_okounkov_dh.tex lines 618--1196
% SHA-256: 51dbbec49fb41b736f3e7d3cd2f2d3bdb8165094a176be237412f435672f5fd5
\section{Mixed volumes of partial Okounkov bodies}
\label{sec:ok-mixed-volumes}
\index{Okounkov body!mixed volume}\index{movable product}

Everything so far attached one convex body to one current and matched its
Euclidean volume with the volume of its \(\mathcal I\)-model envelope.  We
now prove the corresponding mixed formula.  Throughout this section \(X\)
is a connected smooth projective variety over \(\mathbf C\) of dimension
\(n\), and \(\nu\) ranges over the surjective valuations
\[
 \nu:\mathbf C(X)^\times\longrightarrow\mathbf Z^n_{\mathrm{lex}}
 \quad\text{trivial on }\mathbf C^\times ,
\tag{55.26}\label{eq:ok-mvob-valuations}
\]
where \(\mathbf Z^n_{\mathrm{lex}}\) carries the lexicographic order.  Flag
valuations form a subfamily, but every supremum in this section is over the
whole family in \eqref{eq:ok-mvob-valuations}.

\begin{lemma}\label{lem:ok-mvob-valuative-leaves}
Let \(\nu\) be as in \eqref{eq:ok-mvob-valuations}.  Its residue field is
\(\mathbf C\), every nonzero leaf of its valuation filtration is
one-dimensional, and for every nonzero finite-dimensional
\(W\subset\mathbf C(X)\),
\[
                         \#\nu(W^\times)=\dim_{\mathbf C}W.
\tag{55.27}\label{eq:ok-mvob-leaf-count}
\]
The valuation has a center on every proper model of \(\mathbf C(X)\).
Changing a rational trivialization of a line bundle translates its
normalized partial body and hence leaves all mixed volumes unchanged.
\end{lemma}

\begin{proof}
The rational rank of \(\nu\) is \(n\).  Abhyankar's inequality therefore
forces the residue transcendence degree to be zero.  The residue field is
an algebraic extension of the algebraically closed field \(\mathbf C\), so
it is \(\mathbf C\).  After division by any element of value \(a\), the
nonzero leaf
\(\mathcal F_{\succeq a}/\mathcal F_{\succ a}\) embeds in the residue
field.  It is consequently one-dimensional.
Filtering \(W\) by its finitely many values and summing the dimensions of
the leaves proves \eqref{eq:ok-mvob-leaf-count}.  The existence of a center
is the valuative criterion of properness.

If rational frames \(e,e'\) satisfy \(e'=ge\), the functions representing
sections of \(L^{\otimes k}\) differ by multiplication by \(g^{-k}\).
Their normalized valuation sets therefore differ by the fixed translation
\(-\nu(g)\).  Translation invariance of mixed volume proves the last
assertion.
\end{proof}

\begin{remark}[Two normalizations that must be converted, not assumed]
\label{rem:ok-mvob-normalizations}
Two conversions are needed before the source statements can be read against
the conventions of this book.

\emph{Mixed volume.}  For convex bodies \(K_1,\dots,K_n\subset\mathbf R^n\)
this book writes \(V(K_1,\dots,K_n)\) for the coefficient of
\(t_1\cdots t_n\) normalized by \(V(K,\dots,K)=\lambda_n(K)\), as in
\Cref{thm:appD-minkowski-polynomial}.  The source writes \(\vol\) for the
same quantity.  We use \(V\).

\emph{Volume of a metric.}  The source's volume of a psh metric is the one
detected by multiplier ideals, that is
\(\vol(\theta,\varphi)=\int_X(\theta+\ddc P_\theta[\varphi]_{\cI})^n\) of
\Cref{def:complex-current-volume}, and not the bare non-pluripolar mass
\(M_\theta(\varphi)\) of \eqref{eq:complex-npp-mass}.  By
\eqref{eq:complex-mass-volume} one has
\(M_\theta(\varphi)\leq\vol(\theta,\varphi)\), with equality for analytic
singularities and, more generally, for \(\mathcal I\)-model potentials.
The hypotheses of \Cref{def:ok-algebraic-partial-body} and
\Cref{thm:ok-transcendental-package} are therefore exactly the
positive-volume hypothesis appearing in \eqref{eq:ok-mvob-metrics}.  The
mixed quantity \(\vol(\ddc h_1,\dots,\ddc h_n)\) is the mixed
non-pluripolar mass of the \(\cI\)-model envelopes,
\[
 \vol(\ddc h_1,\dots,\ddc h_n)
 =\int_X\bigl(\theta_1+\ddc P_{\theta_1}[\varphi_1]_{\cI}\bigr)
        \wedge\cdots\wedge
        \bigl(\theta_n+\ddc P_{\theta_n}[\varphi_n]_{\cI}\bigr),
\tag{55.28}\label{eq:ok-mvob-mixed-mass}
\]
where \(h_i=h_{0,i}\rme^{-\varphi_i}\) as in
\eqref{eq:ok-partial-metric} and the wedge is the non-pluripolar product
\eqref{eq:complex-global-npp}.
\end{remark}

\subsection{Order and the algebraic upper bound}

\begin{proposition}
\label{prop:ok-mvob-superadditivity}
Let \(L,L'\) be big line bundles on \(X\), let \(h,h'\) be psh metrics on
\(L,L'\), and let \(\nu\) satisfy \eqref{eq:ok-mvob-valuations}.  Assume
that all the partial bodies below are defined, that is, that the relevant
potentials have positive multiplier-ideal volume.  Then
\[
 \Delta_\nu(L,h)+\Delta_\nu(L',h')
   \subseteq\Delta_\nu(L\otimes L',h\otimes h'),
\tag{55.29}\label{eq:ok-mvob-superadditivity}
\]
and, for one fixed smooth background \(\theta\) and
\(\varphi,\psi\in\PSH(X,\theta)\),
\[
 \varphi\preceq_{\cI}\psi
   \quad\Longrightarrow\quad
 \Delta_\nu(\theta,\varphi)\subseteq\Delta_\nu(\theta,\psi).
\tag{55.30}\label{eq:ok-mvob-i-monotone}
\]
\end{proposition}

\begin{proof}
The general-valuation Minkowski inclusion is
\cite[Proposition~10.3.10]{XiaBook}; it gives
\eqref{eq:ok-mvob-superadditivity}.  A naive levelwise product argument
would not suffice, since subadditivity gives
\(\cI(k(\varphi+\varphi'))
 \subseteq\cI(k\varphi)\cI(k\varphi')\), the opposite inclusion.
The implication \eqref{eq:ok-mvob-i-monotone} follows directly from
\eqref{eq:ok-i-order} of \Cref{thm:ok-algebraic-volume-order}.
\end{proof}

\begin{lemma}\label{lem:ok-mvob-finite-level}
Let \(F/\mathbf C\) be finitely generated of transcendence degree \(n\),
let \(X_0\) be an integral projective model of \(F\), and let \(\nu\) be as
in \eqref{eq:ok-mvob-valuations}.  Let
\(0\ne W_i\subset F\) be finite-dimensional, let \(Y\) be the closure of
the image of the joint evaluation map
\[
 X_0\dashrightarrow\prod_{i=1}^n\mathbf P(W_i^\vee),
\]
and suppose that this map is birational onto \(Y\).  If
\[
 \operatorname{span}_{\mathbf Z}
 \{\nu(f)-\nu(g):f,g\in W_i^\times,\ 1\leq i\leq n\}
 =\mathbf Z^n,
\tag{55.31}\label{eq:ok-mvob-finite-level-degree}
\]
then, with \(P_i=\operatorname{conv}\nu(W_i^\times)\) and with \(H_i\)
the hyperplane class from the \(i\)-th factor,
\[
 n!V(P_1,\ldots,P_n)\leq(H_1\cdots H_n\cdot Y).
\tag{55.32}\label{eq:ok-mvob-admissible}
\]
\end{lemma}

\begin{proof}
By \Cref{lem:ok-mvob-valuative-leaves}, choose a basis \(f_{ij}\) of each
\(W_i\), one vector for each value \(a_{ij}=\nu(f_{ij})\).  Choose
\(u_1,\ldots,u_n\in F^\times\) with \(\nu(u_\ell)=e_\ell\).  The residue
field calculation in that lemma gives a graded isomorphism
\[
 \mathbf C[\mathbf Z^n]\simeq\operatorname{gr}_\nu F,
 \qquad \chi^a\longmapsto\operatorname{in}_\nu(u^a).
\]
After rescaling the \(f_{ij}\), their initial forms are \(\chi^{a_{ij}}\).

Let \(S=\mathbf C[x_{ij}]\), multigraded by \(\deg x_{ij}=e_i\), and let
\(I\) be the multihomogeneous prime ideal of \(Y\).  Give \(x_{ij}\) the
lexicographic matrix weight \(a_{ij}\), set
\(Q=\operatorname{in}_{\mathsf A}(I)\), and put
\[
 J=\ker\!\left(S\longrightarrow
 \mathbf C[\mathbf Z^n][t_1,\ldots,t_n],\quad
 x_{ij}\longmapsto\chi^{a_{ij}}t_i\right).
\]
Taking the least-valuation part of each relation among the \(f_{ij}\) gives
only the inclusion
\[
                              Q\subseteq J.
\tag{55.33}\label{eq:ok-mvob-initial-inclusion}
\]
By \cite[Proposition~8.10]{KM19}, a rational scalar weight realizes the
matrix initial ideal; clear its denominators.  The resulting flat
multigraded degeneration preserves every multihomogeneous piece:
\[
 \dim(S/I)_{\mathbf d}=\dim(S/Q)_{\mathbf d}
 \qquad(\mathbf d\in\mathbf N^n).
\tag{55.34}\label{eq:ok-mvob-hilbert-functions}
\]

The ideal \(J\) defines the closure \(Z\) of the monomial map
\[
 (\mathbf C^\times)^n\longrightarrow
 \prod_i\mathbf P(W_i^\vee),\qquad
 z\longmapsto([z^{a_{i0}}:\cdots:z^{a_{ir_i}}])_i.
\]
Condition \eqref{eq:ok-mvob-finite-level-degree} says that the exponent
differences generate the full character lattice.  Thus this map is a
closed immersion on the torus and \(\dim Z=n\).  By
\eqref{eq:ok-mvob-initial-inclusion}, \(Z\) lies in
\(Y_0=\operatorname{MultiProj}(S/Q)\).  Passing to
\(\operatorname{MultiProj}\) removes only multigraded irrelevant torsion;
such torsion vanishes in all sufficiently large multidegrees, and hence
does not change the eventual Hilbert polynomial.  Equation
\eqref{eq:ok-mvob-hilbert-functions} therefore gives
\(\dim Y_0=\dim Y=n\), so \(Z\) is a top-dimensional component of \(Y_0\).
Writing the top-dimensional cycle as
\([Y_0]_n=\sum_Wm_W[W]\), positivity of the multiplicities and nefness of
the globally generated \(H_i\) give
\[
 (H_1\cdots H_n\cdot Y)
 =(H_1\cdots H_n\cdot[Y_0]_n)
 \geq(H_1\cdots H_n\cdot Z).
\tag{55.35}\label{eq:ok-mvob-toric-component}
\]

If some \(W_i\) is one-dimensional, then \(P_i\) is a point, so the mixed
volume in \eqref{eq:ok-mvob-admissible} is zero and the result follows.
Otherwise, for a fixed point of
\(Z\setminus(\mathbf C^\times)^n\), requiring one hyperplane in each factor
to pass through it has total codimension \(n\).  Since the boundary has
dimension at most \(n-1\), the incidence correspondence shows that general
hyperplanes have no common boundary point.  Their intersection is computed
on the dense torus.  The exponent differences generate \(\mathbf Z^n\), so
the Bernstein--Kushnirenko formula \cite{Ber75} has no lattice-index factor
and yields
\[
                    (H_1\cdots H_n\cdot Z)=n!V(P_1,\ldots,P_n).
\]
Together with \eqref{eq:ok-mvob-toric-component}, this proves
\eqref{eq:ok-mvob-admissible}.
\end{proof}

\begin{proposition}\label{prop:ok-mvob-movable-upper}
For big \(\mathbf Q\)-divisors \(D_1,\ldots,D_n\) and every valuation in
\eqref{eq:ok-mvob-valuations},
\[
 n!V(\Delta_\nu(D_1),\ldots,\Delta_\nu(D_n))
 \leq\langle D_1,\ldots,D_n\rangle.
\tag{55.36}\label{eq:ok-mvob-movable-upper}
\]
\end{proposition}

\begin{proof}
By homogeneity, clear denominators.  For
\(W_{i,k}=H^0(X,\mathcal O_X(kD_i))\), use rational trivializations at the
center of \(\nu\); by \Cref{lem:ok-mvob-valuative-leaves}, this only
translates normalized bodies.  For a sufficiently divisible \(p\), every
\(|pD_i|\) is birational.  To obtain the lattice assertion at one common
degree, use the degree-zero fraction field of the section ring of \(D_1\),
which is \(\mathbf C(X)\).  Choose rational functions \(g_j\) with
\(\nu(g_j)=e_j\), and write \(g_j=a_j/b_j\) with \(a_j,b_j\) homogeneous
sections of the same degree.  After passing to one common degree \(p_0\), put
\(s_0=\prod_jb_j\) and
\(s_j=a_j\prod_{\ell\ne j}b_\ell\).  Then all \(s_j\) lie in one
\(W_{1,p_0}\) and \(\nu(s_j)-\nu(s_0)=e_j\).  Replace the previously chosen
birational degree by a common multiple \(p=rp_0\), and replace these sections by
\(s_0^r,s_js_0^{r-1}\).  For \(k=mp\), the sections
\(s_0^{rm},s_js_0^{rm-1}\) retain the unit differences, while the
symmetric-power subsystem induces the Veronese of each degree-\(p\) map.
Thus \Cref{lem:ok-mvob-finite-level} applies at every such \(k\).

Resolve all degree-\(k\) base ideals on one model and write
\[
 \mu_k^*(kD_i)=M_{i,k}+F_{i,k},\qquad
 M_{i,k}\ \text{basepoint free},\quad F_{i,k}\geq0.
\]
The projection formula, in the form of \Cref{thm:appF-projection}, gives
\[
 n!V(\operatorname{conv}\nu(W_{1,k}^\times),\ldots,
      \operatorname{conv}\nu(W_{n,k}^\times))
 \leq M_{1,k}\cdots M_{n,k}.
\]
The nef approximants \(k^{-1}M_{i,k}\) lie below the pullbacks of \(D_i\).
The nef-approximation inequality
\cite[Propositions~2.12 and~2.13]{BFJ09} therefore gives
\[
 M_{1,k}\cdots M_{n,k}
 \leq k^n\langle D_1,\ldots,D_n\rangle.
\]
Divide by \(k^n\).  The normalized finite-level convex hulls converge in
Hausdorff distance by \cite[Proposition~10.3.4]{XiaBook}; continuity of
mixed volume in \Cref{thm:appD-minkowski-polynomial} proves
\eqref{eq:ok-mvob-movable-upper}.
\end{proof}

\subsection{Uniformity in the valuation}

\begin{lemma}\label{lem:ok-mvob-potential-sandwich}
Let \(u_j,u\in\PSH(X,\theta)_{>0}\) be normalized model potentials with
\(u_j\to u\) in \(d_S\).  After passing to a subsequence there are
increasing normalized model potentials \(\eta_j\) and decreasing normalized
model potentials \(\psi_j\) such that
\[
 \eta_j\leq u_j\leq\psi_j,\qquad
 \eta_j\leq u\leq\psi_j,\qquad
 \eta_j,\psi_j\longrightarrow u\quad\text{in }d_S.
\tag{55.37}\label{eq:ok-mvob-potential-sandwich}
\]
For finitely many sequences, the same subsequence may be chosen for all of
them.
\end{lemma}

\begin{proof}
The sandwich theorem \cite[Proposition~6.2.3]{XiaBook} supplies the first
pair of inequalities and the two convergences.  If its upper sequence is
not model, replace it by \(P_\theta[\psi_j]\).  Monotonicity preserves
decrease and the upper inequality, while model projection preserves the
\(d_S\)-limit.

Fix \(j\).  For \(k\geq j\),
\(\eta_j\leq\eta_k\leq u_k\).  Closure of the \(P\)-preorder under
simultaneous \(d_S\)-limits \cite[Proposition~6.2.5]{XiaBook}, applied to
the constant sequence \(\eta_j\) and the tail \((u_k)_{k\geq j}\), gives
\(\eta_j\preceq_Pu\).  Since both potentials are normalized and model,
this is \(\eta_j\leq u\).  Similarly,
\(u_k\leq\psi_k\leq\psi_j\) for \(k\geq j\), so order closure and modelness
give \(u\leq\psi_j\).  This proves the missing two inequalities.  For
finitely many sequences, apply the construction successively and take the
finite diagonal subsequence.
\end{proof}

\begin{lemma}\label{lem:ok-mvob-uniform-sandwich}
For \(1\leq i\leq n\), suppose that nonempty compact convex sets satisfy
all four inclusions
\[
 A_{i,j}\subseteq C_{i,j}\subseteq B_{i,j},\qquad
 A_{i,j}\subseteq D_{i,j}\subseteq B_{i,j}.
\tag{55.38}\label{eq:ok-mvob-convex-sandwich}
\]
Assume, for constants \(v,M>0\) independent of \(i,j\), that
\[
 \lambda_n(A_{i,j})\geq v,\qquad
 \lambda_n(B_{1,j}+\cdots+B_{n,j})\leq M,\qquad
 \lambda_n(B_{i,j})-\lambda_n(A_{i,j})\longrightarrow0.
\]
Then
\[
 \left|V(C_{1,j},\ldots,C_{n,j})
       -V(D_{1,j},\ldots,D_{n,j})\right|\longrightarrow0.
\]
\end{lemma}

\begin{proof}
For each \(i,j\), translate the whole quadruple so that
\(0\in A_{i,j}\).  Every body is then contained in
\(S_j=B_{1,j}+\cdots+B_{n,j}\), and
\(v\leq\lambda_n(S_j)\leq M\).  Send the John ellipsoid of \(S_j\) to the
unit ball by an affine map with linear part \(L_j\).  John's theorem
\cite[Theorem~10.12.2]{Sch14} gives
\[
 \mathbf B^n\subseteq L_jS_j+b_j\subseteq n\mathbf B^n,\qquad
 \frac{\lambda_n(\mathbf B^n)}{M}\leq|\det L_j|
 \leq\frac{n^n\lambda_n(\mathbf B^n)}{v}.
\]
All transformed bodies lie in one compact set, their inner volumes stay
positive, and their inner--outer volume gaps tend to zero.  From any
subsequence, Blaschke selection \cite[Theorem~1.8.7]{Sch14} gives a further
subsequence on which all
\(4n\) bodies converge.  At the limit, the four inclusions, equality of the
positive inner and outer volumes, and strict volume monotonicity under a
proper inclusion of full-dimensional convex bodies imply
\(\widehat A_i=\widehat B_i=\widehat C_i=\widehat D_i\).  Continuity of
mixed volume gives the conclusion after transformation.  The lower bound
for \(|\det L_j|\) transfers it back.  Since every subsequence has such a
further subsequence, the full sequence converges.
\end{proof}

\begin{proposition}\label{prop:ok-mvob-uniform-ds}
For \(1\leq i\leq n\), let \(L_i\) be a big line bundle, choose a smooth
form \(\theta_i\in c_1(L_i)\), let
\(\varphi_{i,j},\varphi_i\in\PSH(X,\theta_i)_{>0}\), and assume that
\(\varphi_{i,j}\to\varphi_i\) in \(d_S\).
Then
\[
 \sup_\nu\left|
 V(\Delta_\nu(\theta_1,\varphi_{1,j}),\ldots,
   \Delta_\nu(\theta_n,\varphi_{n,j}))
 -V(\Delta_\nu(\theta_1,\varphi_1),\ldots,
   \Delta_\nu(\theta_n,\varphi_n))\right|\longrightarrow0.
\tag{55.39}\label{eq:ok-mvob-uniform-ds}
\]
\end{proposition}

\begin{proof}
Replace all potentials by normalized model envelopes.  If
\eqref{eq:ok-mvob-uniform-ds} failed, an approximate maximizer of the
supremum would give, after a subsequence, valuations \(\nu_j\) and
\(\delta>0\) for which the displayed difference is at least \(\delta\).
Use the common subsequence in \Cref{lem:ok-mvob-potential-sandwich} and set
\[
 \begin{aligned}
 A_{i,j}&=\Delta_{\nu_j}(\theta_i,\eta_{i,j}),&
 B_{i,j}&=\Delta_{\nu_j}(\theta_i,\psi_{i,j}),\\
 C_{i,j}&=\Delta_{\nu_j}(\theta_i,\varphi_{i,j}),&
 D_{i,j}&=\Delta_{\nu_j}(\theta_i,\varphi_i).
 \end{aligned}
\]
Discard a finite initial segment.  Continuity of the diagonal masses and
positivity of the limits then put every one of these bodies on the
positive-volume locus.  The two sandwiches in
\eqref{eq:ok-mvob-potential-sandwich} and
\eqref{eq:ok-mvob-i-monotone} give all four inclusions in
\eqref{eq:ok-mvob-convex-sandwich}.  The diagonal volume formula and
\(d_S\)-continuity give
\[
 \lambda_n(B_{i,j})-\lambda_n(A_{i,j})\longrightarrow0,\qquad
 \lambda_n(A_{i,j})\longrightarrow
 \frac1{n!}\vol(\theta_i,\varphi_i)>0,
\]
independently of \(\nu_j\).  Moreover,
\[
 B_{1,j}+\cdots+B_{n,j}
 \subseteq\Delta_{\nu_j}\!\left(\sum_i\theta_i,
                                  \sum_i\psi_{i,j}\right)
 \subseteq\Delta_{\nu_j}\!\left(\bigotimes_iL_i\right)
\]
by \eqref{eq:ok-mvob-superadditivity} and containment in the complete
series.  The last body has Euclidean volume
\(\vol(\bigotimes_iL_i)/n!\), independent of \(j\) and \(\nu_j\).
\Cref{lem:ok-mvob-uniform-sandwich} now contradicts \(\delta\).  This
proves \eqref{eq:ok-mvob-uniform-ds}.  Notice that the conclusion concerns
mixed-volume values uniformly in \(\nu\), not Hausdorff convergence
uniformly in \(\nu\).
\end{proof}

\subsection{Rational analytic singularities}

\begin{definition}\label{def:ok-rational-analytic}
A psh metric has \emph{rational analytic singularities} if locally its
weight is
\[
 c\log\!\left(\sum_j|f_j|^2\right)+O(1),\qquad c\in\mathbf Q_{>0}.
\]
This is the rational subcase of
\Cref{def:complex-analytic-singularities}; the rationality is needed for
the divisor and section arguments below.
\end{definition}

\begin{lemma}\label{lem:ok-mvob-effective-shift}
Let \(A,B\) be big \(\mathbf Q\)-divisors, let \(E\geq0\) be a
\(\mathbf Q\)-divisor, and suppose
\(A\sim_{\mathbf Q}aB+bE\) with \(a,b\in\mathbf Q_{\geq0}\).  Then
\[
 a\Delta_\nu(B)+b\nu(E)\subseteq\Delta_\nu(A).
\tag{55.40}\label{eq:ok-mvob-effective-shift}
\]
\end{lemma}

\begin{proof}
Clear the denominators in the divisor relation and in \(a,b\).  Multiplying
a section of the corresponding multiple of \(B\) by the required power of
the canonical section of \(E\) produces a section of the corresponding
multiple of \(A\).  Valuation vectors add under multiplication.  Divide by
the common multiple and take the closed convex hull to obtain
\eqref{eq:ok-mvob-effective-shift}.
\end{proof}

\begin{proposition}\label{prop:ok-mvob-analytic-case}
Assume that the psh metrics \(h_i\) on the big line bundles \(L_i\) have
rational analytic singularities and that their curvature currents are
K\"ahler currents.  Then \eqref{eq:ok-mvob-metrics} holds.
\end{proposition}

\begin{proof}
On a common projective log resolution \(\pi:Y\to X\), write
\[
 \pi^*c_1(L_i,h_i)=[D_i]+R_i,\qquad
 A_i:=\pi^*L_i-D_i\in\operatorname{Pic}(Y)_{\mathbf Q},
\]
where \(D_i\geq0\) is rational and \(R_i\) has bounded potentials.  The
class \(\alpha_i=c_1(A_i)=\{R_i\}\) is nef and big.  Birational invariance
of the non-pluripolar product gives
\[
 \int_X\bigwedge_i c_1(L_i,h_i)=\alpha_1\cdots\alpha_n.
\tag{55.41}\label{eq:ok-mvob-resolution-mass}
\]
For every valuation of the common function field, the log-resolution
formula \cite[Remark~10.3.1, equation~(10.22)]{XiaBook} gives
\[
 \Delta_\nu(L_i,h_i)=\Delta_\nu(A_i)+\nu(D_i).
\tag{55.42}\label{eq:ok-mvob-resolution-body}
\]
Translations do not affect mixed volume.  Applying
\Cref{prop:ok-mvob-movable-upper}, and using nefness, yields
\[
 n!V(\Delta_\nu(L_1,h_1),\ldots,\Delta_\nu(L_n,h_n))
 \leq\alpha_1\cdots\alpha_n.
\tag{55.43}\label{eq:ok-mvob-analytic-upper}
\]

For the reverse inequality, fix an ample \(\mathbf Q\)-divisor \(H\) on
\(Y\).  Choose one rational \(c>0\) such that every \(cA_i-H\) is big, and
choose \(E_i\geq0\) with \(E_i\sim_{\mathbf Q}cA_i-H\).  For rational
\(\varepsilon>0\), set \(B_i=A_i+\varepsilon H\).  Then \(B_i\) is ample
and
\[
 A_i\sim_{\mathbf Q}
 \frac1{1+c\varepsilon}B_i
 +\frac{\varepsilon}{1+c\varepsilon}E_i.
\tag{55.44}\label{eq:ok-mvob-kodaira}
\]
By \cite[Theorem~1.1]{Wil25}, there is an admissible flag, allowed to depend
on \(\varepsilon\), whose valuation \(\nu_\varepsilon\) satisfies
\[
 n!V(\Delta_{\nu_\varepsilon}(B_1),\ldots,
      \Delta_{\nu_\varepsilon}(B_n))=B_1\cdots B_n.
\tag{55.45}\label{eq:ok-mvob-wilms}
\]
Apply \Cref{lem:ok-mvob-effective-shift} to
\eqref{eq:ok-mvob-kodaira}.  Translation invariance and
\eqref{eq:ok-mvob-wilms} give
\[
 \sup_\nu V(\Delta_\nu(A_1),\ldots,\Delta_\nu(A_n))
 \geq\frac{B_1\cdots B_n}{n!(1+c\varepsilon)^n}.
\tag{55.46}\label{eq:ok-mvob-analytic-lower}
\]
Let rational \(\varepsilon\downarrow0\).  Equations
\eqref{eq:ok-mvob-resolution-mass}--\eqref{eq:ok-mvob-analytic-lower}
prove the proposition.  No maximizing valuation was used: at each
\(\varepsilon\) a possibly different flag supplies the lower bound.
\end{proof}

\subsection{The positive-volume formula}

\begin{lemma}\label{lem:ok-mvob-i-model-approximation}
Let \(h=h_0\rme^{-\varphi}\) be a psh metric of positive
multiplier-ideal volume and put \(u=P_\theta[\varphi]_{\cI}\).  Then
\(u\sim_{\cI}\varphi\), the partial bodies of \(u\) and \(\varphi\) agree,
\(u\) is a model potential with
\(M_\theta(u)=\vol(\theta,\varphi)>0\), and there are rational-analytic
potentials \(u_j\) such that
\[
 u_j\longrightarrow u\quad\text{in }d_S,\qquad
 \theta+\ddc u_j\ \text{is a K\"ahler current}.
\tag{55.47}\label{eq:ok-mvob-approximating}
\]
\end{lemma}

\begin{proof}
The first two assertions follow from
\Cref{thm:complex-i-projection} and the \(\cI\)-invariance of partial
bodies.  It is also model: the model envelope
\(P_\theta[u]\) is \(\cI\)-equivalent to \(u\), so maximality of the
\(\cI\)-envelope gives \(P_\theta[u]\leq u\), while normalized \(u\) is a
competitor for its model envelope and gives the reverse inequality.  The
mass--volume identity for an \(\cI\)-model potential is
\eqref{eq:complex-mass-volume}.  Thus the assumed ideal volume, rather than
the possibly smaller raw mass of \(\varphi\), is the positive mass of
\(u\).  The potential \(u\) is \(\cI\)-good, and the rational-analytic
K\"ahler-current approximation in \cite[Theorem~7.1.1]{XiaBook} gives
\eqref{eq:ok-mvob-approximating}.
\end{proof}

\begin{theorem}\label{thm:ok-mvob-metrics}
Let \(L_i\) be big line bundles on \(X\), and let \(h_i\) be psh metrics
with \(\vol(\theta_i,\varphi_i)>0\), in the normalization of
\Cref{rem:ok-mvob-normalizations}.  Then
\[
 \vol(\ddc h_1,\dots,\ddc h_n)
 =n!\,\sup_\nu V\!\left(
    \Delta_\nu(L_1,h_1),\dots,\Delta_\nu(L_n,h_n)\right).
\tag{55.48}\label{eq:ok-mvob-metrics}
\]
\end{theorem}

\begin{proof}
For each \(i\), apply \Cref{lem:ok-mvob-i-model-approximation} and write
\(u_i=P_{\theta_i}[\varphi_i]_{\cI}\).  The replacement preserves every
body, and by definition \eqref{eq:ok-mvob-mixed-mass} the left side is the
mixed non-pluripolar mass of the \(u_i\).  Choose the approximants
\(u_{i,j}\) in \eqref{eq:ok-mvob-approximating}, using one common index for
the finitely many \(i\).  At every \(j\),
\Cref{prop:ok-mvob-analytic-case} gives
\[
 \int_X\bigwedge_i(\theta_i+\ddc u_{i,j})
 =n!\sup_\nu V(\Delta_\nu(\theta_1,u_{1,j}),\ldots,
                    \Delta_\nu(\theta_n,u_{n,j})).
\]
The mixed non-pluripolar masses on the left converge by
\cite[Theorem~6.2.1]{XiaBook}.  The right side converges by
\Cref{prop:ok-mvob-uniform-ds}; explicitly, for the corresponding functions
\(f_j,f\) of \(\nu\),
\[
 \left|\sup_\nu f_j(\nu)-\sup_\nu f(\nu)\right|
 \leq\sup_\nu|f_j(\nu)-f(\nu)|\longrightarrow0.
\]
Taking the limit proves \eqref{eq:ok-mvob-metrics}.
\end{proof}
% END EXACT COMPONENT CH55_METRIC_ONLY_CHAIN

% --- Primary-source attestations (packet-only; not part of live Chapter 55) ---
\begin{quote}
\textbf{XiaBook primary-source attestation.}
The persistent file \texttt{tmp/pdfs/SGPT\_final\_20260718\_primary.pdf}
has SHA-256
\texttt{c816eea5f06822021fa8dd75b70f82724ee47e6dcd6830729f5ecc109296b706}.
Direct inspection confirms: Proposition~6.2.5 closes the \(P\)-preorder
under simultaneous \(d_S\)-limits; Proposition~10.3.4 gives normalized
finite-level Hausdorff convergence; Proposition~10.3.10 gives the
Minkowski inclusion for partial Okounkov bodies; and Remark~10.3.1,
equation~(10.22), gives the log-resolution translation formula.
These are printed on book pages 225, 335, 340, and 338, respectively
(physical PDF pages 241, 351, 356, and 354).

\textbf{Schneider locator attestation.}
In the 2014 second expanded edition of \emph{Convex Bodies: The
Brunn--Minkowski Theory}, Theorem~10.12.2 on page 588 gives the John
ellipsoid normalization used above, and Theorem~1.8.7 on page 63 is the
Blaschke selection theorem.  Theorem~1.8.6 only states closedness of the
space of convex bodies and is not the cited selection result.
\end{quote}

% BEGIN EXACT COMPONENT CH55_LINE_BUNDLE_TARGET
% Source: manuscript/parts/part5/ch55_okounkov_dh.tex lines 1198--1217
% SHA-256: 87ffb791c287e180631231c97b2435c037e7573adf51174192efbb47979d3a25
\begin{theorem}\label{thm:ok-mvob-line-bundles}
For big line bundles \(L_1,\ldots,L_n\) on \(X\),
\[
 \langle L_1,\dots,L_n\rangle
 =n!\,\sup_\nu V\!\left(
     \Delta_\nu(L_1),\dots,\Delta_\nu(L_n)\right).
\tag{55.49}\label{eq:ok-mvob-line-bundles}
\]
\end{theorem}

\begin{proof}
Choose psh metrics with minimal singularities in the classes \(c_1(L_i)\).
Their partial bodies are the ordinary Okounkov bodies by
\Cref{thm:ok-algebraic-volume-order}.  The mixed non-pluripolar mass of
minimal currents is the positive intersection number by
\cite[Definition~1.17]{BEGZ10}; its projective numerical interpretation is
the movable product of \cite[Propositions~2.12--2.13]{BFJ09}.
Thus \Cref{thm:ok-mvob-metrics} specializes exactly to
\eqref{eq:ok-mvob-line-bundles}.
\end{proof}
% END EXACT COMPONENT CH55_LINE_BUNDLE_TARGET
